\section{関連研究}
\subsection{物体認識モデル}
YOLO\cite{ref_YOLO}は画像を格子に分割し，各格子からボックスとクラスを1回の推論で予測する．
本稿ではYOLOの系列のYOLO11x\cite{ref_YOLO11}を用いる．
DETR\cite{ref_DETR}はTransformerによる集合予測として物体認識を定式化し，重複するボックスの除去を行わずにボックスとクラスを出力する．
いずれもCOCO\cite{ref_COCO}の80クラスで学習されたモデルが公開されており，学習時のクラス以外を出力することはできない．

\subsection{文章・画像埋め込みモデル}
CLIP\cite{ref_CLIP}は画像と文章の対の対照学習により，画像エンコーダとテキストエンコーダを同じ空間に埋め込む．
SigLIP\cite{ref_SigLIP}は，対照学習の目的関数を，候補全体で正規化するソフトマックスから，画像と文章の対ごとのシグモイドに置き換えたモデルである．
いずれも，クラス名を「a photo of a 〇〇」のような文章に埋め込み，画像との類似度を比較することで，追加の学習なしに任意のクラスへの分類（ゼロショット分類）を行うことができる．
OWL-ViT\cite{ref_OWLViT}やGrounding DINO\cite{ref_GroundingDINO}は，このような埋め込みを物体認識に組み込み，テキストで指定したクラスの物体を検出する．
本稿は，物体認識モデルと文章・画像埋め込みモデルを組み合わせ，ボックスの検出は物体認識モデルで1回だけ行い，クラスの推定を文章・画像埋め込みモデルで繰り返す点で，これらのモデルと異なる．

\subsection{WordNet}
WordNet\cite{ref_WordNet}は英語の名詞を同義語の集合（Synset）にまとめ，Synsetの間の上位語・下位語の関係を記述した語彙データベースである．
例えばdogのSynsetの下位語にはworking dogやhunting dog等が含まれ，下位語を順にたどることで細かい粒度の単語が得られる．
名詞のSynsetは，最上位のカテゴリ（unique beginner）を根とする木構造をなしている．

\subsection{階層的な分類と確信度}
確信度が低い場合に粗いクラスを出力する考え方は，画像分類において正解率と細かさのトレードオフとして定式化されている\cite{ref_Hedging}．
Dengらは，クラスの階層上で正解率の下限を満たしつつ情報量が最大となるクラスを選ぶ方法を提案した．
また，CHiLS\cite{ref_CHiLS}は，ゼロショット分類において，各クラスの下位のクラスの名前を文章・画像埋め込みモデルに与え，その確率を上位のクラスへ集約することで正解率を改善した．
これらはいずれも画像全体のクラスが事前に定めた集合に含まれる分類の問題であり，本稿は，検出を伴い，学習時に無い物体に7806語のWordNetの語彙から粒度を選んで命名する点で異なる．
出力の確率を正解率と一致させる校正には，ロジスティック回帰による方法\cite{ref_Platt}や温度によるスケーリング\cite{ref_Calibration}があり，校正の程度は期待校正誤差で評価される\cite{ref_Calibration}．
